\section{Spy vs Spy}

\begin{quotex}
Un bel giorno ... un bel giorno per morire
\flright{\textsc{Giacomo Puccini}\footnote{\url{https://www.youtube.com/watch?v=RnCmeNa5vXs}}}
\end{quotex}


\paragraph{The Americans}

The fourth season of the FX series \emph{The Americans} begins next week. The setting is suburban Washington, DC, in the early 1980s. The premise is that Elizabeth and Philip Jennings, outwardly the parents of a typical American family, are deeply embedded KGB agents. Raising a family and running a travel agency would be overwhelming for most people. Yet they have the time and energy to do the spy stuff, too.

They are nearly superhuman with mad martial arts skills, both armed and unarmed. That should require frequent training at the gym and practice at the shooting range, but we never see that. They are masters of disguise and social engineering, bending people to their will. Apparently, they have a license to kill, which they often resort to. Yet they manage to remain sympathetic, since the American public is accustomed to killing as long as it is done for a higher cause. For example, the Godfather movies are also like that: murder is not ``personal'', since such motives would make it immoral. There are also real life examples that I'm sure you can find.

The Jennings' are noble since they fight for their fatherland and the communist ideology. In comparison, the American FBI agents, always close on their tail but not quite there, are lost. For example, Stan the FBI man, seems to work out of a sense of duty, or more likely to get out of the house, but is unclear about any higher purpose. After all, the USA line is that ``they hate us for our freedom'' and they come here ``for all our stuff''.

Yet for all his ``freedom'', Stan is emotionally crippled and spiritually empty. His family is disintegrating and he and his wife fumble around with various New Age cults (like est) for life's purpose. Elizabeth Jennings, on the other hand, has no doubts about her life purpose. She even retains a somewhat warrior ascetic attitude in the midst of the consumer society. When her husband tells her how ``easy'' it is to live in USA, she responds, ``It is different, but not better.''

Curiously, KGB agents were active, behind the scenes, in liberal political causes. In particular, we can identify these left wing causes:

\begin{itemize}


\item Civil rights movements: they were active in a cell

\item Iran Contra affair: they supported the Sandanistas in Nicaragua

\item Nuclear disarmament: they encouraged their teenaged daughter to participate in demonstrations. (The daughter, unaware of her parents' true identities, joined a liberal church.)

\item The World Council of Churches was infiltrated by communists
\end{itemize}


So ``Russian'' interference in American politics was, in the 70s and 80s, more than just Internet ``misinformation''.

\paragraph{Secret Agent Man}


A totally different point of view can be found in the TV series from the 60s, Danger Man UK) or Secret Agent\footnote{\url{https://www.google.com/url?sa=t\&rct=j\&q=\&esrc=s\&source=web\&cd=8\&cad=rja\&uact=8\&ved=0ahUKEwjTsYqR2rLLAhVFbR4KHclDCGYQtwIISTAH\&url=https\%3A\%2F\%2Fwww.youtube.com\%2Fwatch\%3Fv\%3D6iaR3WO71j4\&usg=AFQjCNFq2P3q43phFIlIY085y1isGKaU4g\&sig2=g0cendtLTM1734oOsfRrgg\&bvm=bv.116274245,d.dmo}} (USA). John Drake was quite sure of his mission. Moreover, he, too, relied on social engineering to solve his cases with minimal use of force. Watch some episodes if you want to see what life in a ``white republic'' was like. Or even better, watch \emph{Bye, Bye, Birdie}. When my millennial aged son was a young teen, I called him to the TV to see the big dance scene in the town center. I asked him if anything seemed odd about it. He caught on right away!

\paragraph{White Supremacy}


The elections in the USA have brought attention to ``white supremacy''. It is hard to see this as much of a factor, since it involves little more than a few web sites. Nevertheless, it is the ``other'' that the left loves to hate.

Nevertheless, I did a little research and biological racism seems entrenched in certain circles, usually associated with anti-Judaism and anti-Christianity. This is far from a reactionary position. In his critique of Fascism, Evola reminds us:

\begin{quotex}
In truth, personality and liberty can be conceived only on the basis of the \emph{liberation of the individual from naturalistic, biological, and primitively individualistic bonds} that characterize the pre-state and pre-political forms in a purely social and utilitarian-contractual sense.
\end{quotex}


They seem to be attracted to neo-Darwinism, as sort of a ``tough minded'' worldview (in William James' sense). Unfortunately, the so-called ``neo reactionaries'' are not immune from this attitude. Another Hermetist, Valentin Tomberg, in his commentary on the Arcana of the Moon, explains the danger of this attitude:

\begin{quotex}
What is the state of intelligence which has abandoned all metaphysics and has decided to limit itself solely to ``objective facts of the senses''? What is most characteristic of this state is that
\emph{intelligence no longer moves forwards but backwards}. It looks to

the least developed and the most primitive for the cause and explanation of what is most developed and most advanced in the process of evolution.
\emph{Thus, it looks for the efficient cause of the world not in the heights of creative consciousness but rather in the depths of the unconscious,} and instead of advancing and elevating itself towards God, it retreats into matter.
\end{quotex}


Evola regards the same movement as ``categorical'', or a false feeling of self-transcendence, which

\begin{quotex}
can give to the individual the illusory, momentary sensation of an exalted, intense life, moreover such sensations conditioned by a regression, by a diminution of the personality and true liberty.
\end{quotex}


\paragraph{Putrefaction of the Modern Soul}


This regression is evidenced by what such beings simply cannot see. Specifically, they fail to see what Christianity had brought to Europe. One particular web site claimed that Christianity is like a disease that takes time to manifest itself; therefore the supremacy of Christian Europe is an illusion. Then, even though a de-Christianized Europe seems to be the problem, he claims that Europe is still infected. So he has a ``heads I win, tails you lose'' attitude if you want to debate him.

On the other hand, Evola could see Medieval Europe as an ``\emph{heroic, differentiated, and spiritual civilization}''. Anyone who cannot recognize that is hardly likely to recognize it today. There are two errors in biological racism. The first is that genetics is not constant, not just due to mixing, but also to the accumulation of mutations over generations. A materially prosperous civilization is more likely to be declining genetically.

The second is that it ignores the race of the soul and psychical heredity, which is even more important. It is not an ideology or religion or worldview that kills the soul, but rather that a weak soul will embrace a disordered worldview. Perhaps such men are reacting to the putrefaction of the modern soul, and then confuse the cause and the effect.

Europe is just a land mass. The true destiny of its people is once again to be Holy and Roman and an Empire.

\paragraph{Catholic vs Orthodox}


This question keeps arising in certain circles. Since Gornahoor is focused on the esoteric teachings, he seldom concerns himself with the exoteric aspects. The follow up question is why didn't Rene Guenon and Julius Evola ever comment on the Eastern churches?

The answer is that there is no difference between the two from the esoteric perspective. They are, or at least used to be, in a schismatic relationship to each other (the mutual excommunications of 1054 were lifted some 20 years ago). Now a schism is primarily an organizational dispute, unlike a heresy which involves a difference in doctrine. Too often the focus is on accidental and historically contingent differences. However, there are no essential differences.

This is shown by the examples of \textbf{St Nicodemus} and \textbf{St Theophan the Recluse}, who aggregated the \emph{Philokalia} for the Greeks and Russians respectively. Nevertheless, they could recognize In \emph{Spiritual Combat} by \textbf{Lorenzo Scupoli} a work of deep spirituality. They translated that book into their languages. Moreover, St Nicodemus translated the \emph{Spiritual Exercises} of \textbf{St Ignatius Loyola} into Greek, calling it simply the work of a ``wise man''. Any other attitude does not serve a good purpose.

\paragraph{Sticks and Stones}


\begin{quotex}
Sticks and stones may break my bones, but names will never harm me.
\end{quotex}


That was the guiding philosophy of my youth, so the schoolyard bantering in the Republican debates does not concern me. However, it upsets many people and is, in any case, inappropriate in such a venue. Nevertheless, the opposition does use names with the intent to harm people. In particular, there is vitriol against \textbf{Donald Trump}, and by extension his supporters, using words like bigoted, racist, sexist, deplorable, and so on. These words are intended to put someone ``beyond the pale'', yet the pundit class does not deem them inappropriate. I heard the editor of a political journal blame ``racism and nativism'' as the root cause for so many things she didn't like; she may as well have tried to explain fire by phlogiston.

Speaking of ``stones'', the wannabe adult on stage, \textbf{John Kasich}, was vividly threatening to go to war with Russia to defend Finland, inter alia, against an imminent invasion. While that bellicose attitude was astonishing enough, what was even more astonishing was how
\emph{normal} it sounded to everyone. I don't believe it has ever been brought up since.


\flrightit{Posted on 2016-03-08 by Cologero}

\begin{center}* * *\end{center}

\begin{footnotesize}
\begin{sffamily}
\texttt{Harun on 2016-03-09 at 06:25 said:}

Most of the points pseudo-reactionary make are based on flawed modern premises, like a purely materialistic world view.

They explain the current state of Europe to race, and thus darwinism and evolutionism. They are effectively eliminating any possibility of solving a spiritual problem.

Also they fail to realize biological race is a modern concern, not very relevant in the middle age or in the cosmopolitic roman empire.


\hfill

\texttt{Alistair Fraser on 2016-03-11 at 03:32 said:}

the host has made a superb job on this gornahoor site of clearly distinguishing the phenomena of race distinctions/types and their mis-interpreted connotations into modern consciousness of which many people seem unable to comprehend due to the mass mediated agitating version that is recycled in modern times with a purpose mostly to hostility.

I thought this quote from evola from 1971 was a sort of parallel mirror with that race question, this time on the mis-perceptions of ``magic'' and the accusations that can be thrown from that
... ... ... ... ... ...

``JE: As I formerly said, there are people who want to create myths, even of the most comical sort, in order to find self-satisfaction. It was even stated in France that in my presence Black Masses were celebrated with young blondes every week. Given my condition we could recall a German saying that can be translated as ``he got it wrong because it was true.'' Therefore I can only say in this regard that during the years nineteen-twenty-seven, twenty-eight and twenty-nine, I organised what is called the Ur Group, which did not deal with magic, but rather with initiatic operations. And after this period it is important to distinguish between the magical practices and the initiatic ones which are quite different -- and I must say that from the Ur Group arose a work of three volumes called Introduction to Magic which is composed of anonymous monographies since the names that appear are not those of the authors of the diverse subjects, translations and comments. Thus, I believe that anyone who aims to explore such a subject seriously, and in a multifaceted way, should resort to this indispensable work. A third edition in Italian appeared with the same title.''

\end{sffamily}
\end{footnotesize}

